\writeSectionFrame{\presentationTitle}{Beispiel aus der realen Welt}
\begin{frame}{Dekorator wird eingesetzt z.B. in}
	\begin{itemize}
 		\item Java I/O
        \item Java Collections Framework
 		\item JakartaEE
 		\item Spring
        \item GUI-Frameworks
 	\end{itemize}
\end{frame}
\note{
    Dekorator ist ein Muster, das häufig in Frameworks und Softwareprojekten eingesetzt wird, besonders in der Java-Welt. 
}

\begin{frame}{Java I/O}
	\begin{itemize}
 		\item Kernkomponenten: InputStream, OutputStream, Reader, Writer
        \item Konkrete Komponenten: FileInputStream, FileOutputStream, BufferedReader, BufferedWriter, etc.
 		\item Dekoratoren: Klassen wie BufferedInputStream, DataInputStream, GZIPOutputStream, PrintWriter
 	\end{itemize}
\end{frame}
\note{
    In Java I/O werden die Dekoratoren verwendet, um zusätzliche Funktionalitäten wie Pufferung, Datenumwandlung, Komprimierung usw. einzubringen. 
}

\begin{frame}{Decorator}
	\begin{center}
 		\includegraphics[scale=0.35]{\BILDERPFAD/realwelt1.png}
 	\end{center}
\end{frame}

\note{
	Wir möchten einen neuen Stream programmieren, den LowerCaseStream, der den Inhalt in Kleinbuchstaben umwandelt. 
}

\begin{frame}{Input-Stream-Hierarchie}
\begin{tikzpicture}[scale=0.8, transform shape]
    % Definiere eine Klasse
    \umlclass[x=5, y=0]{InputStream}{
    }{
    }

    \umlclass[x=0, y=0]{FileInputStream}{
    }{
    }

    \umlclass[x=10, y=0]{PipedInputStream}{
    }{
    }

    \umlclass[x=5, y=-2.5]{FilterInputStream}{
    }{
    }

    \umlclass[x=5, y=-5]{BufferedInputStream}{
    }{
    }

    \umlclass[x=10, y=-5]{DataInputStream}{
    }{
    }

    \umlclass[x=0, y=-5]{PushbackInputStream}{
    }{
    }
    
    \umlinherit{FileInputStream}{InputStream}
    \umlinherit{PipedInputStream}{InputStream}
    \umlinherit{FilterInputStream}{InputStream}
    \umlinherit{BufferedInputStream}{FilterInputStream}
    \umlinherit{DataInputStream}{FilterInputStream}
    \umlinherit{PushbackInputStream}{FilterInputStream}
\end{tikzpicture}
\end{frame}

\begin{frame}{Output-Stream-Hierarchie}
\begin{tikzpicture}[scale=0.8, transform shape]
    % Definiere eine Klasse
    \umlclass[x=5, y=0]{OutputStream}{
    }{
    }

    \umlclass[x=0, y=0]{FileOutputStream}{
    }{
    }

    \umlclass[x=10, y=0]{PipedOutputStream}{
    }{
    }

    \umlclass[x=5, y=-2.5]{FilterOutputStream}{
    }{
    }

    \umlclass[x=5, y=-5]{BufferedOutputStream}{
    }{
    }

    \umlclass[x=10, y=-5]{DataOutputStream}{
    }{
    }

    \umlclass[x=0, y=-5]{PrintStream}{
    }{
    }
    
    \umlinherit{FileOutputStream}{OutputStream}
    \umlinherit{PipedOutputStream}{OutputStream}
    \umlinherit{FilterOutputStream}{OutputStream}
    \umlinherit{BufferedOutputStream}{FilterOutputStream}
    \umlinherit{DataOutputStream}{FilterOutputStream}
    \umlinherit{PrintStream}{FilterOutputStream}
\end{tikzpicture}
\end{frame}

\begin{frame}[fragile]{In eine Datei schreiben}
	\begin{exampleblock}{LowerCaseInputStreamDialog.java}
		\begin{lstlisting}
private void initializeFile() {
    File file = new File(FILENAME);
    try (FileWriter fileWriter = new FileWriter(file);
         BufferedWriter bufferedWriter = 
            new BufferedWriter(fileWriter)) {
        bufferedWriter.write("....\n");
        bufferedWriter.write("...");
    } catch (IOException e) {
        e.printStackTrace();
    }
}
 		\end{lstlisting}
	 \end{exampleblock}
\end{frame}

\begin{frame}{Unser neue Dekorator}
\begin{tikzpicture}[scale=0.8, transform shape]
    % Definiere eine Klasse
    \umlclass[x=5, y=0]{InputStream}{
    }{
    }

    \umlclass[x=0, y=0]{FileInputStream}{
    }{
    }

    \umlclass[x=10, y=0]{PipedInputStream}{
    }{
    }

    \umlclass[x=5, y=-2.5]{FilterInputStream}{
    }{
    }

    \umlclass[x=5, y=-5]{BufferedInputStream}{
    }{
    }

    \umlclass[x=10, y=-5]{DataInputStream}{
    }{
    }

    \umlclass[x=0, y=-5]{PushbackInputStream}{
    }{
    }

    \umlclass[x=10, y=-2.5]{LowerCaseInputStream}{
    }{
    }
    
    \umlinherit{FileInputStream}{InputStream}
    \umlinherit{PipedInputStream}{InputStream}
    \umlinherit{FilterInputStream}{InputStream}
    \umlinherit{BufferedInputStream}{FilterInputStream}
    \umlinherit{DataInputStream}{FilterInputStream}
    \umlinherit{PushbackInputStream}{FilterInputStream}
    \umlinherit{LowerCaseInputStream}{FilterInputStream}
\end{tikzpicture}
\end{frame}

\begin{frame}[fragile]{Unser neuer Dekorierer}
	\begin{exampleblock}{LowerCaseInputStream.java}
		\begin{lstlisting}
public class LowerCaseInputStream extends FilterInputStream {
	public LowerCaseInputStream(InputStream in) {
		super(in);
	}
	
	public int read() throws IOException {
		int c = super.read();
		if (c == -1) {
			return c; 
		} else {
			return Character.toLowerCase((char)c);
		}
	}
}
		\end{lstlisting}
	 \end{exampleblock}
\end{frame}

\begin{frame}[fragile]{Dialogmethode für den Test}
	\begin{exampleblock}{MainDialog.java}
		\begin{lstlisting}
private void start() {
    int c;
    
    try {
        InputStream in = new LowerCaseInputStream(
                new BufferedInputStream(
                    new FileInputStream(FILENAME)));
        while ((c = in.read()) >= 0) {
            System.out.print((char) c);
        }
        in.close();
    } catch (IOException e) {
        e.printStackTrace();
    }
}
        \end{lstlisting}
	\end{exampleblock}
\end{frame}

\begin{frame}[fragile]{Java Collections Framework}
	\begin{exampleblock}{Beispiel}
		\begin{lstlisting}
List<String> list = new ArrayList<>();
List<String> synchronizedList = Collections.
    synchronizedList(list);
List<String> synchronizedUnmodifiableList = 
    Collections.unmodifiableList(synchronizedList);

        \end{lstlisting}
	\end{exampleblock}
\end{frame}
\note{
    Im Java Collection Framework kann man den Dekorator bspw. bei der Klasse Collections beobachten. Hier wird eine Liste in eine synchronisierte und unmodifizierbare Liste verpackt, wobei das ursprüngliche Verhalten der Liste beibehalten wird.
}

\begin{frame}[fragile]{Servlet-Filter in Jakarta EE}
	\begin{exampleblock}{Beispiel}
		\begin{lstlisting}
public class CompressionFilter implements Filter {
    @Override
    public void doFilter(ServletRequest request, 
                ServletResponse response, FilterChain chain) 
            throws IOException, ServletException {
        HttpServletResponse httpResp = (HttpServletResponse) 
            response;
        GZIPResponseWrapper wrappedResponse = 
            new GZIPResponseWrapper(httpResp);
        chain.doFilter(request, wrappedResponse);
        wrappedResponse.finishResponse();
    }
}
        \end{lstlisting}
	\end{exampleblock}
\end{frame}
\note{
    In Webanwendungen, die mit Java EE (jetzt Jakarta EE) entwickelt wurden, werden Servlet-Filter verwendet, um Anfragen und Antworten zu modifizieren. Kernkomponenten: HttpServletRequest, HttpServletResponse, Dekoratoren sind Filter wie CompressionFilter, LoggingFilter, AuthenticationFilter, die zusätzliche Funktionen wie Komprimierung, Protokollierung und Authentifizierung hinzufügen. Ein CompressionFilter komprimiert eine Antwort, bevor sie an den Client gesendet wird.
}
